
\section{Probabilistic Tape Diagrams}\label{sec:tapediagrams}
%
Recall that a \emph{monoidal signature} is a tuple $(\sort, \sign, \ari, \coar)$ where $\sort$ is a set of basic sorts, hereafter denoted by $A,B,\dots$, $\sign$ is a set of generators, denoted by $s,t, \dots$, and $\ari,\coar \colon \sign \to \sort^*$ assign to each symbol its arity and coarity (words over $\sort$). From a monoidal signature $(\sort, \sign, \ari, \coar)$, one can freely generate the strict symmetric monoidal category $\DiagS$:  objects are words in $\sort^*$; arrows are \emph{string diagrams} \cite{joyal1991geometry,selinger2010survey}. These can be regarded  as the terms generated by the following grammar (where $A,B \in \sort$ and $s \in \Sigma$)
\begin{equation}\label{stringdiagramGrammar}
\setlength{\arraycolsep}{3pt} % spaziatura orizzontale
\renewcommand{\arraystretch}{1.1} % opzionale: un po' di respiro verticale
\begin{array}{rcccccccccccccccc}
c & ::= & \id{A} & \mid & \id{\uno} & \mid & \gen & \mid & \symmt{A}{B} & \mid & c ; c & \mid & c \per c
\end{array}
\end{equation}
modulo the axioms of strict symmetric monoidal categories. A \emph{monoidal theory} $\mathbb{T}=(\sign, E)$ consists of a monoidal signature equipped with a set $E$ of pairs of arrows of $\DiagS$ with same source and target. Let $=_{\mathbb{T}}$ be an equivalence relation on arrows  of $\DiagS$ obtained as the congruence closure (w.r.t. $;$ and $\otimes$) of $E$. We write $\DiagT$ for the category obtained as the quotient of $\DiagS$ by $=_{\mathbb{T}}$.

\medskip

The category $\CatT{\DiagT}$, obtained by applying the construction from Section~\ref{sec:syntactic} to $\DiagT$, coincides by definition with the category of \emph{tape diagrams} from \cite{bonchi2025tapediagramsmonoidalmonads}. Objects of $\CatT{\DiagT}$ are elements of $(\sort^*)^*$ which we often write as \emph{polynomials} $P=\Piu[i=1][n]{\Per[j=1][m_i]{A_{i,j}}}$. 
We will call \emph{monomials} of $P$ the $n$ words $\Per[j=1][m_i]{A_{i,j}}$. For instance, the monomials of $(A \per B) \piu \uno$ are $A \per B$ and $1$. We denote monomials by $U,V,\dots$

Arrows of $\CatT{\DiagT}$ enjoy an intuitive diagrammatic representation specified by the following two layered grammar.
\begin{equation*}\label{tapesDiagGrammar} 
    \setlength{\tabcolsep}{2pt}
    \begin{tabular}{rc c@{$\,\mid\,$}c@{$\,\mid\,$}c@{$\,\mid\,$}c@{$\,\mid\,$}c@{$\,\mid\,$}c@{$\,\mid\,$}c@{$\,\mid\,$}c@{$\,\mid\,$}c@{$\,\mid\,$}c}
        $c$  & $\Coloneqq$ &  $\wire{A}$ & $ \tikzfig{empty} $ & $ \Cgen{\gen}{A}{B}  $ & $ \Csymm{A}{B} $ & $ \tikzfig{seq_compC}   $ & $  \tikzfig{par_compC}$ \\
        $\t$ & $\Coloneqq$ & $\tikzfig{/tapes/cipriano/pcomonoid}$ & $\Tcounit{U}$  & $ \Tcirc{c}{U}{V}$ & $\Tunit{U}$  & $\Tmonoid{U}$ \\ 
        && $\Twire{U}$ & $ \tikzfig{empty} $    & $ \Tsymmp{U}{V} $ & $ \tikzfig{tapes/seq_comp}  $ & $  \tikzfig{tapes/par_comp}$  
    \end{tabular}
\end{equation*} 
The first layer corresponds to \eqref{stringdiagramGrammar} while the second to \eqref{tapesGrammar}; diagrams from the first layer are string diagrams, those from the second are called tapes. Note that (a) string diagram can occur inside tapes, (b) string diagrams have type $U\to V$, for $U,V\in \sort^*$ and vertical composition corresponds to $\per$, (c) tapes have type $P\to Q$ for $P,Q\in (\sort^*)^*$ and vertical composition corresponds to $\piu$. 

The identity $\id\zero$ is rendered as the empty tape $\tikzfig{empty}$, while $\id\uno$ is $\tikzfig{tapes/empty}$: a tape filled with the empty string diagram. 
For a monomial $U \!=\! A_1\dots A_n$, $\id U$ is depicted as a tape containing  $n$ wires labelled by $A_i$. For instance, $\id{AB}$ is rendered as $\TRwire{A}{B}$. When clear from the context, we will simply represent it as a single wire  $\Twire{U}$ with the appropriate label.
Similarly, for a polynomial $P = \Piu[i=1][n]{U_i}$, $\id{P}$ is obtained as a vertical composition of tapes, as illustrated below on the left. 
%
\begin{equation}\label{ex:tape}
\scalebox{0.82}{$ 
    \id{AB \piu \uno \piu C} = \!\!\!\begin{aligned}\begin{gathered} \TRwire{A}{B} \\[-1.8mm] \Twire{\uno} \\[-1.8mm] \Twire{C} \end{gathered}\end{aligned}
    \quad
    \codiag{A\piu B \piu C} = \!\!\tikzfig{tapes/examples/codiagApBpC}
    \quad
    \cobang{AB \piu B \piu C} = \!\tikzfig{tapes/examples/cobangABpBpC}
    \quad
    \tikzfig{tapes/cipriano/algt}
$} \end{equation}
%
\noindent The codiagonal $\codiag{U} \colon  U \piu U \!\to\! U$ is represented as a merging of tapes, $\diagp{U}\colon U\to U\oplus U $ as a splitting of tapes labeled with $p\in(0,1)$, the cobang $\cobang{U} \colon \zero \!\to\! U$ is a tape closed on its left boundary, while $\bang{U} \colon U \to \zero $ is closed on the right.
Exploiting the definitions in \eqref{eq:indmoncopca}, we can construct $\codiag{P},\diagp{P},\cobang{P},\bang{P}$ for arbitrary polynomials.
For example, $\codiag{A\piu B \piu C}$ and $\cobang{AB \piu B \piu C}$ are depicted as the second and third diagrams above. The last diagram is the tape for 
$\diagp{A};(\diagq{A}\oplus \id{A}) ; (\bang{A} \oplus \codiag{A})$. 
For an arbitrary $\t \colon P \to Q$, we write $\Tbox{\t}{P}{Q}$. 

The graphical representation embodies several axioms such as \eqref{eq:TapeFunctAxioms} and those of monoidal categories. 
%The remaining axioms are shown in Figure~\ref{fig:tapesax}. %\marginpar{Mancano gli assiomi dei tape e delle simmetrie del per dentro i 
Those axioms that are not implicit in the graphical representation are shown in Figure~\ref{fig:tapesax}. %\marginpar{Mancano gli assiomi dei tape e delle simmetrie del per dentro i tape... No?} 
%
%
%

\medskip

By Corollary~\ref{cor:isotapematrices}, we know that $\CatT{\DiagT}$ is isomorphic to $\stmat{\DiagT^+}$: tapes represent exactly stochastic matrices of subprobability distributions of string diagrams in $\DiagT$. Figure \ref{fig:dictionary} illustrates  
how the isomorphism translates tape diagrams into such matrices. For instance, the rightmost diagram in \eqref{ex:tape} corresponds to the $1\times 1$ matrix $\begin{pmatrix} (1-p)(1-q) \cdot \wire{A}\end{pmatrix}$ while the following tape diagram on the left represents the $2\times 2$ stochastic matrix on the right.
\[
\resizebox{0.95\linewidth}{!}{$
\tikzfig{tapes/cipriano/esempiomatrice}
\quad
\text{\raisebox{1.7ex}{$
\begin{pNiceMatrix}[first-col,first-row]
	\rotatebox{90}{$\Lsh$} & AB & C \\
	DE & p\cdot \tikzfig{par_compCMAT} & q\cdot \Cgen{h}{}{} \\
	F & (1-p)\cdot \star & (1-q)\cdot(\tikzfig{seq_compCMAT1} +_r \tikzfig{seq_compCMAT2})
\end{pNiceMatrix}
$}}
$}
\]

\begin{comment}

\begin{figure}[t]
\centering
\setlength{\tabcolsep}{2pt} % riduce lo spazio orizzontale
\renewcommand{\arraystretch}{0.9} % riduce lo spazio verticale
\resizebox{0.92\linewidth}{!}{% scala globale per stare nei margini
\begin{tabular}{r@{\;}c@{\;}l @{\qquad} r@{\;}c@{\;}l @{\qquad} r@{\;}c@{\;}l}
    % 1st row
    \tikzfig{tapes/ax/symminv_left} &$\stackrel{(\symmp\text{-inv})}{=}$& \tikzfig{tapes/ax/symminv_right}  
    & 
    \tikzfig{tapes/ax/symmnat_left}  &$\stackrel{(\symmp\text{-nat})}{=}$&  \tikzfig{tapes/ax/symmnat_right}
    & 
    \tikzfig{cb/symm_inv_left} & $\axeq{\sigma\text{-inv}}$ & \tikzfig{cb/symm_inv_right} 
    \\
    % 2nd row
    \tikzfig{tapes/whiskered_ax/monoid_assoc_left} &$\stackrel{(\codiag{}\text{-as})}{=}$& \tikzfig{tapes/whiskered_ax/monoid_assoc_right}  
    &
    \tikzfig{tapes/whiskered_ax/monoid_unit_left} &$\stackrel{(\codiag{}\text{-un})}{=}$& \Twire{U}
    &
    \tikzfig{tapes/whiskered_ax/monoid_comm_left} &$\stackrel{(\codiag{}\text{-sym})}{=}$& \Tmonoid{U}
    \\
    % 3rd row
    \tikzfig{tapes/cipriano/monoid_assoc_left} & $\stackrel{(\diagp{}{}\text{-as})}{=}$ & \tikzfig{tapes/cipriano/monoid_assoc_right} 
    &
    \scalebox{0.8}{\tikzfig{tapes/cipriano/idempotency}}&$\stackrel{(\codiag{}\text{-idem})}{=}$&   \Twire{U}
    &
    \tikzfig{tapes/cipriano/monoid_comm_left} & $\stackrel{(\diagp{}{}\text{-sym})}{=}$ & \tikzfig{tapes/cipriano/pcomonoidbar}
    \\
    % Final row (scaled to match others)
    \multicolumn{9}{c}{%
        \scalebox{0.85}{%
        \begin{tabular}{r@{\;}c@{\;}l @{\quad} r@{\;}c@{\;}l @{\quad} r@{\;}c@{\;}l @{\quad} r@{\;}c@{\;}l}
            \tikzfig{tapes/cipriano/pcanatleft} & $\stackrel{(\diagp{}\text{-nat})}{=}$ & \tikzfig{tapes/cipriano/pcanatright} 
            &
            \tikzfig{tapes/cipriano/monoidnatnewright} & $\stackrel{(\codiag{}\text{-nat})}{=}$ & \tikzfig{tapes/cipriano/monoidnatnewleft} 
            &
            \tikzfig{tapes/cipriano/bangnatleft} & $\stackrel{(\,\bangp{}\text{-nat})}{=}$ & \tikzfig{tapes/cipriano/bangnatright} 
            &
            \tikzfig{tapes/cipriano/cobangnatleft} & $\stackrel{(\cobang{}\text{-nat})}{=}$ & \tikzfig{tapes/cipriano/cobangnatright}
        \end{tabular}
        }% fine scalebox
    }
\end{tabular}
}% fine resizebox
\caption{Axioms for probabilistic tape diagrams.}
\label{fig:tapesax}
\end{figure}
\end{comment}




\begin{figure}[t]
\centering
\setlength{\tabcolsep}{2pt}
\renewcommand{\arraystretch}{0.9}
\resizebox{0.92\linewidth}{!}{%
\begin{tabular}{r@{\;}c@{\;}l @{\qquad} r@{\;}c@{\;}l @{\qquad} r@{\;}c@{\;}l}
    % 1st row
    \mylabelbis{ax:tapes:symminv}{\ensuremath{\symmp\text{-inv}}}%
    \tikzfig{tapes/ax/symminv_left} &$\stackrel{\eqref{ax:tapes:symminv}}{=}$& \tikzfig{tapes/ax/symminv_right}
    &
    \mylabelbis{ax:tapes:symmnat}{\ensuremath{\symmp\text{-nat}}}%
    \tikzfig{tapes/ax/symmnat_left} &$\stackrel{\eqref{ax:tapes:symmnat}}{=}$& \tikzfig{tapes/ax/symmnat_right}
    &
    \mylabelbis{ax:tapes:sigmainv}{\ensuremath{\sigma\text{-inv}}}%
    \tikzfig{cb/symm_inv_left} &$\stackrel{\eqref{ax:tapes:sigmainv}}{=}$& \tikzfig{cb/symm_inv_right}
    \\
    % 2nd row
    \mylabelbis{ax:tapes:codiagas}{\ensuremath{\codiag{}\text{-as}}}%
    \tikzfig{tapes/whiskered_ax/monoid_assoc_left} &$\stackrel{\eqref{ax:tapes:codiagas}}{=}$& \tikzfig{tapes/whiskered_ax/monoid_assoc_right}
    &
    \mylabelbis{ax:tapes:codiagun}{\ensuremath{\codiag{}\text{-un}}}%
    \tikzfig{tapes/whiskered_ax/monoid_unit_left} &$\stackrel{\eqref{ax:tapes:codiagun}}{=}$& \Twire{U}
    &
    \mylabelbis{ax:tapes:codiagsym}{\ensuremath{\codiag{}\text{-sym}}}%
    \tikzfig{tapes/whiskered_ax/monoid_comm_left} &$\stackrel{\eqref{ax:tapes:codiagsym}}{=}$& \Tmonoid{U}
    \\
    % 3rd row
    \mylabelbis{ax:tapes:diagpas}{\ensuremath{\diagp{}{}\text{-as}}}%
    \tikzfig{tapes/cipriano/monoid_assoc_left} &$\stackrel{\eqref{ax:tapes:diagpas}}{=}$& \tikzfig{tapes/cipriano/monoid_assoc_right}
    &
    \mylabelbis{ax:tapes:diagpidem}{\ensuremath{\diagp{}\text{-idem}}}%
    \scalebox{0.8}{\tikzfig{tapes/cipriano/idempotency}} &$\stackrel{\eqref{ax:tapes:diagpidem}}{=}$& \Twire{U}
    &
    \mylabelbis{ax:tapes:diagpsym}{\ensuremath{\diagp{}{}\text{-sym}}}%
    \tikzfig{tapes/cipriano/monoid_comm_left} &$\stackrel{\eqref{ax:tapes:diagpsym}}{=}$& \tikzfig{tapes/cipriano/pcomonoidbar}
    \\
    % Final row
    \multicolumn{9}{c}{%
        \scalebox{0.85}{%
        \begin{tabular}{r@{\;}c@{\;}l @{\quad} r@{\;}c@{\;}l @{\quad} r@{\;}c@{\;}l @{\quad} r@{\;}c@{\;}l}
            \mylabelbis{ax:tapes:diagpnat}{\ensuremath{\diagp{}\text{-nat}}}%
            \tikzfig{tapes/cipriano/pcanatleft} &$\stackrel{\eqref{ax:tapes:diagpnat}}{=}$& \tikzfig{tapes/cipriano/pcanatright}
            &
            \mylabelbis{ax:tapes:codiagnat}{\ensuremath{\codiag{}\text{-nat}}}%
            \tikzfig{tapes/cipriano/monoidnatnewright} &$\stackrel{\eqref{ax:tapes:codiagnat}}{=}$& \tikzfig{tapes/cipriano/monoidnatnewleft}
            &
            \mylabelbis{ax:tapes:bangnat}{\ensuremath{\,\bangp{}\text{-nat}}}%
            \tikzfig{tapes/cipriano/bangnatleft} &$\stackrel{\eqref{ax:tapes:bangnat}}{=}$& \tikzfig{tapes/cipriano/bangnatright}
            &
            \mylabelbis{ax:tapes:cobangnat}{\ensuremath{\,\cobang{}\text{-nat}}}%
            \tikzfig{tapes/cipriano/cobangnatleft} &$\stackrel{\eqref{ax:tapes:cobangnat}}{=}$& \tikzfig{tapes/cipriano/cobangnatright}
        \end{tabular}
        }%
    }
\end{tabular}
}%
\caption{Axioms for probabilistic tape diagrams.}
\label{fig:tapesax}
\end{figure}
%
%



\medskip 
It is now crucial to recall from \cite{bonchi2025tapediagramsmonoidalmonads} that $\per$ can be defined not only on string diagrams but also on tapes: for $\t_1 \colon P \to Q$, $\t_2 \colon R \to S$,  $\t_1 \per \t_2 \defeq \LW{P}{\t_2} ; \RW{S}{\t_1} $ where $\LW{P}{\cdot}$, $\RW{S}{\cdot}$ are the left and right whiskerings, defined as in Table~\ref{tab:producttape}. For objects $P = \Piu[i]{U_i}$ and $Q = \Piu[j]{V_j}$, 
\begin{equation}\label{def:productPolynomials} 
P \per Q \defeq \Piu[i]{\Piu[j]{U_iV_j}}\text{.}
\end{equation}

Theorem 27 in \cite{bonchi2025tapediagramsmonoidalmonads} guarantees that $(\CatT{\DiagT} ,\otimes, \uno)$ is a symmetric monoidal category. Such a category is monoidally enriched over $\Cat{PCA}$, i.e., for all $p\in(0,1)$ and properly typed $\t,\s_1,\s_2$:
\begin{equation}\label{eq:monoidalenrichment} \t \otimes (\s_1+_p \s_2)= (\t \otimes \s_1)+_p(\t \otimes \s_2) \qquad  (\s_1+_p \s_2)  \otimes \t= ( \s_1 \otimes \t )+_p(\s_2 \otimes  \t) \qquad \star \otimes \t= \star = \t\otimes \star \end{equation}
 Most importantly, the two monoidal categories $(\CatT{\DiagT}, \piu, \zero)$ and $(\CatT{\DiagT}, \otimes, \uno)$ interact through the laws of (right strict) rig categories (see \cite{laplaza_coherence_1972,johnson2021bimonoidal}).

\begin{theorem}[From \cite{bonchi2025tapediagramsmonoidalmonads}]
$(\,\CatT{\DiagT}, \oplus, \otimes, \zero, \uno \,)$ is a right strict rig category.
\end{theorem}

\begin{table}[t]
    \begin{center}
    {
        \hfill {\tiny
  \[\begin{array}{c}
  \toprule
        \def\arraystretch{1.2}
        \begin{array}{cc}
            \begin{array}{@{}l}
                % \toprule
                \dl{P}{Q}{R} \colon P \per (Q\piu R)  \to (P \per Q) \piu (P\per R) \vphantom{\symmt{P}{Q}} \\
                \midrule
                \dl{\zero}{Q}{R} \defeq \id{\zero} \vphantom{\symmt{P}{\zero} \defeq \id{\zero}} \qquad
                \dl{U \piu P'}{Q}{R} \defeq (\id{U\per (Q \piu R)} \piu \dl{P'}{Q}{R});(\id{U\per Q} \piu \symmp{U\per R}{P'\per Q} \piu \id{P'\per R}) \vphantom{\symmt{P}{V \piu Q'} \defeq \dl{P}{V}{Q'} ; (\Piu[i]{\tapesymm{U_i}{V}} \piu \symmt{P}{Q'})} \\
                % \bottomrule
            \end{array}
            &
            \begin{array}{@{}l}
                % \toprule
                \symmt{P}{Q} \colon P\per Q \to Q \per P, \text{ with } P = \Piu[i]{U_i} \\
                \midrule
                \symmt{P}{\zero} \defeq \id{\zero} \qquad
                \symmt{P}{V \piu Q'} \defeq \dl{P}{V}{Q'} ; (\Piu[i]{\tapesymm{U_i}{V}} \piu \symmt{P}{Q'})
                \phantom{\quad}
                % \bottomrule
            \end{array}
        \end{array}\\
        \midrule
        \begin{array}{rclrcl|rclrcl}
            \midrule
            \LW U {\id\zero} &\defeq& \id\zero& \LW U {\t_1 \piu \t_2} &\defeq& \LW U {\t_1} \piu \LW U {\t_2} &  \RW U {\id\zero} &\defeq& \id\zero & \RW U {\t_1 \piu \t_2} &\defeq& \RW U {\t_1} \piu \RW U {\t_2} \\
            \LW U {\tape{c}} &\defeq& \tape{\id U \per c} &    \LW U {\t_1 ; \t_2} &\defeq& \LW U {\t_1} ; \LW U {\t_2}&  \RW U {\tape{c}} &\defeq& \tape{c \per \id U} & \RW U {\t_1 ; \t_2} &\defeq& \RW U {\t_1} ; \RW U {\t_2} \\
            \LW U {\symmp{V}{W}} &\defeq& \symmp{UV}{UW}       &       && & \RW U {\symmp{V}{W}} &\defeq& \symmp{VU}{WU}&  &&  \\
                        \LW U {\diagp V} &\defeq& \diagp{UV} &  \LW U {\bang V} &\defeq& \bang{UV}& \RW U {\diagp V} &\defeq& \diagp{VU} & \RW U {\bang V} &\defeq& \bang{VU} \\
            \LW U {\codiag V} &\defeq& \codiag{UV} &  \LW U {\cobang V} &\defeq& \cobang{UV}& \RW U {\codiag V} &\defeq& \codiag{VU} & \RW U {\cobang V} &\defeq& \cobang{VU} \\
            \hline \hline
            \LW{\zero}{\t} &\defeq& \id{\zero} &  \LW{W\piu S'}{\t} &\defeq& \LW{W}{\t} \piu \LW{S'}{\t} & \RW{\zero}{\t} &\defeq& \id{\zero} &
            \RW{W \piu S'}{\t} &\defeq& \dl{P}{W}{S'} ; (\RW{W}{\t} \piu \RW{S'}{\t}) ; \Idl{Q}{W}{S'} \\
            \hline \hline
            \multicolumn{12}{c}{
                \t_1 \per \t_2 \defeq \LW{P}{\t_2} ; \RW{S}{\t_1}   \quad \text{ ( for }\t_1 \colon P \to Q, \t_2 \colon R \to S   \text{ )}
            }
            \\
            \bottomrule
        \end{array}
    \end{array}        
      \]
      }
      \hfill
      \caption{Inductive definitions of  left distributor $\dl{P}{Q}{R}$ and $\otimes$-symmetry $\symmt{P}{Q}$ (top); monomial and polynomial whiskerings (center); definition of $\per$ (bottom).}\label{tab:producttape}
       }
    \end{center}
\end{table}


By means of the isomorphism in Corollary~\ref{cor:isotapematrices}, we can equip $\stmat{\DiagT^+}$ with an additional product $\otimes$. This is described as follows. On objects $P\otimes Q$ is defined as in \eqref{def:productPolynomials}. On arrows, first define, for all $d_1\in \DiagT^+[U_1,V_1]$ and $d_2\in \DiagT^+[U_2,V_2]$,  $d_1\otimes^+d_2 \in \DiagT^+[U_1U_2,V_1V_2]$ as
\[d_1\otimes^+ d_2 (c) \defeq \sum_{\{(c_1,c_2) \mid c_1\otimes c_2 =c\}} d_1(c_1)\cdot d_2(c_2) \qquad \text{ for }c\in \DiagT[U_1U_2,V_1V_2]\text{.}\]
(recall from Section~\ref{sec:pca} that $ \DiagT^+[U,V]=\Dis(\DiagT[U,V])$). Then for matrices, it is defined as the Kronecker product using $\otimes^+$ as multiplication. 
More precisely, let $M\colon\bigoplus_{i=1}^nU_i \to \bigoplus_{i'=1}^{n'}U_{i'}$ and $N\colon\bigoplus_{j=1}^m V_j \to \bigoplus_{j'=1}^{m'} V_{j'}$, $M\otimes N$ is a matrix of size $n'm'\times nm$ whose rows are indexed by pairs $(i',j')\in \{1,\dots, n'\}\times \{1,\dots m'\}$ and columns by pairs $(i,j)\in \{1,\dots, n\}\times \{1,\dots m\}$. Assuming that $M_{i',i}=p_{i',i} \cdot d_{i',i}$
and $N_{j',j}=q_{j',j} \cdot e_{j',j}$, $M\otimes N$ is defined for all indexes $(i',j'),(i,j)$ as
$(M\otimes N)_{(i',j'),(i,j)} \defeq (p_{i',i}\cdot q_{j',j}) \cdot (d_{i',i}\otimes^+ e_{j',j}) $.
%
\begin{corollary}
$\CatT{\DiagT}$ and $\stmat{\DiagT^+}$ are isomorphic as rig categories.
\end{corollary}




Obviously, the results illustrated in this section also hold for $\CatT{\Diag{\Sigma}}$, namely the category of tapes over string diagrams generated by the monoidal signature $\Sigma$, rather than over the theory $\mathbb{T}$. It is often convenient to use both  $\CatT{\Diag{\Sigma}}$ and  $\CatT{\Diag{\mathbb{T}}}$. These convex biproduct categories are related by the morphism
\[\CatT{Q_\mathbb{T}} \colon \CatT{\Diag{\Sigma}} \to \CatT{\Diag{\mathbb{T}}}\]
obtained by applying the functor $\CatT{-}$ to $Q_\mathbb{T}\colon \Diag{\Sigma} \to \Diag{\mathbb{T}}$, namely the symmetric monoidal functor mapping each string diagram into its $=_{\mathbb{T}}$-equivalence class. The equivalence induced by $\CatT{Q_\mathbb{T}}$ can be easily characterised as the congruence $\tapeeqT$ generated by the following rules.

\begin{equation}\label{eq:congr2}
        \scalebox{0.85}{$
        \centering
        \begin{array}{c}
        \begin{array}{c@{\qquad\qquad}c@{\qquad\qquad}c@{\qquad\qquad}c}
            \inferrule*[right=($\mathbb{T}$)]{c =_{\mathbb{T}} d}{\tape{c} \tapeeqT \tape{d}}
            &
            \inferrule*[right=($\textsc{R}$)]{-}{\t \tapeeqT \t}
            &
            \inferrule*[right=($\textsc{S}$)]{\t_1 \tapeeqT \t_2}{\t_2 \tapeeqT \t_1}
            &    
            \inferrule*[right=($\textsc{T}$)]{\t_1 \tapeeqT \t_2 \quad \t_2 \tapeeqT \t_3}{\t_1 \tapeeqT \t_3}
            % &
            % \inferrule*[right=($E$)]{ (t,t') \in E \and U \in \sort^\star}{ t_U \tapeeqT t'_U}
        \end{array}
        \\[10pt]
        \begin{array}{c@{\qquad}c@{\qquad}c}
            \inferrule*[right=($;$)]{\t_1 \tapeeqT \t_2 \quad \s_1 \tapeeqT \s_2}{\t_1;\s_1 \tapeeqT \t_2;\s_2}
            &
            \inferrule*[right=($\piu$)]{\t_1 \tapeeqT \t_2 \quad \s_1 \tapeeqT \s_2}{\t_1\piu\s_1 \tapeeqT \t_2 \piu \s_2}
            &
            \inferrule*[right=($\per $)]{\t_1 \tapeeqT \t_2 \quad \s_1 \tapeeqT \s_2}{\t_1\per \s_1 \tapeeqT \t_2 \per \s_2}       
        \end{array}
        \end{array}
        $}
\end{equation}

\begin{proposition}\label{prop:quotienttheory}
For all $\s , \t \in \CatT{\Diag{\Sigma}}$, $\s \tapeeqT \t$ iff $\CatT{Q_{\mathbb{T}}}(\s) = \CatT{Q_{\mathbb{T}}}(\t)$.
\end{proposition}


% $    & $  $ & $  $ & $  $  
\begin{figure}
\centering
\begingroup
\setlength{\tabcolsep}{4pt}
\renewcommand{\arraystretch}{1.05}
\small
\begin{tabular}{c c c}
\toprule
Symbol & Diagram & Matrix\\
\midrule

$\diagp{U}$
&
\tikzfig{/tapes/cipriano/pcomonoid}
&
\raisebox{-0.4ex}{$
\begin{pNiceMatrix}[first-col,first-row]
\rotatebox{90}{$\Lsh$} & U \\
U & p\cdot \wire{U} \\
U & (1-p)\cdot \wire{U}
\end{pNiceMatrix}
$}
\\[0.2em]

$\bangp{U}$
&
\Tcounit{U}
&
\raisebox{-0.4ex}{$
\begin{pNiceMatrix}[first-col,first-row]
\rotatebox{90}{$\Lsh$} \\
U
\end{pNiceMatrix}
$}
\\[0.2em]

$\tapeFunct{c}$
&
\Tcirc{c}{U}{V}
&
\raisebox{-0.4ex}{$
\begin{pNiceMatrix}[first-col,first-row]
\rotatebox{90}{$\Lsh$} & U \\
V & 1\cdot \Cgen{c}{U}{V}
\end{pNiceMatrix}
$}
\\[0.2em]

$\cobang{U}$
&
\Tunit{U}
&
\raisebox{-0.4ex}{$
\begin{pNiceMatrix}[first-col,first-row]
\rotatebox{90}{$\Lsh$} & U
\end{pNiceMatrix}
$}
\\[0.2em]

$\codiag{U}$
&
\Tmonoid{U}
&
\raisebox{-0.4ex}{$
\begin{pNiceMatrix}[first-col,first-row]
\rotatebox{90}{$\Lsh$} & U & U \\
U & 1\cdot \wire{U} & 1\cdot \wire{U}
\end{pNiceMatrix}
$}
\\[0.2em]

$\id{U}$
&
\Twire{U}
&
\raisebox{-0.4ex}{$
\begin{pNiceMatrix}[first-col,first-row]
\rotatebox{90}{$\Lsh$} & U \\
U & 1\cdot \wire{U}
\end{pNiceMatrix}
$}
\\[0.2em]

$\id{\zero}$
&
\tikzfig{empty}
&
\raisebox{-0.4ex}{$
\begin{pNiceMatrix}[first-col,first-row]
\rotatebox{90}{$\Lsh$}
\end{pNiceMatrix}
$}
\\[0.2em]

$\sigma_{U,V}^{\piu}$
&
\Tsymmp{U}{V}
&
\raisebox{-0.4ex}{$
\begin{pNiceMatrix}[first-col,first-row]
\rotatebox{90}{$\Lsh$} & U & V \\
V & 0\cdot \star_{U,V} & 1\cdot \wire{V} \\
U & 1\cdot \wire{U} & 0\cdot \star_{V,U}
\end{pNiceMatrix}
$}
\\[0.2em]

$\t ; \t$
&
\tikzfig{tapes/seq_comp}
&
matrix multiplication
\\[0.2em]

$\t \piu \t$
&
\tikzfig{tapes/par_comp}
&
direct sum of matrices
\\

\bottomrule
\end{tabular}
\endgroup
\caption{Dictionary of correspondences}\label{fig:dictionary}
\end{figure}

